\section{Empirical Results}
\begin{frame}[t]{AI and Solution Quality}
\small
\begin{columns}[T,onlytextwidth]
\begin{column}{.60\textwidth}
\figureplaceholder{2}{Average solution quality}{4.7cm}
\end{column}
\begin{column}{0.36\textwidth}
\begin{block}{Finding}
Quality increases by \textbf{0.245 SD} for human-only teams, \textbf{0.373 SD} for individuals with AI, and \textbf{0.392 SD} for teams with AI, relative to solo work without AI.
\end{block}
\begin{alertblock}{Interpretation}
AI-assisted individuals achieve quality comparable to human-only teams. AI can supply some of the performance benefits associated with collaboration.
\end{alertblock}
\end{column}
\end{columns}

\end{frame}

\begin{frame}[t]{Task Familiarity and Expertise}
\small
\begin{columns}[T,onlytextwidth]
\begin{column}{.60\textwidth}
\figureplaceholder{4}{Quality by task familiarity}{4.7cm}
\end{column}
\begin{column}{0.36\textwidth}
\begin{block}{Finding}
Employees with less routine involvement in product development improve when working individually with AI. Their performance approaches teams containing experienced colleagues.
\end{block}
\begin{alertblock}{Interpretation}
AI expands access to task-relevant expertise, helping employees contribute to innovation beyond their routine responsibilities.
\end{alertblock}
\end{column}
\end{columns}

\end{frame}

\begin{frame}[t]{Bridging Functional Boundaries}
\small
\begin{columns}[T,onlytextwidth]
\begin{column}{.60\textwidth}
\figureplaceholder{5}{Commercial and technical orientation}{4.7cm}
\end{column}
\begin{column}{0.36\textwidth}
\begin{block}{Finding}
Without AI, R\&D employees propose more technical ideas and commercial employees emphasize market considerations. With AI, the two distributions become more similar.
\end{block}
\begin{alertblock}{Interpretation}
AI helps employees incorporate perspectives from another function, producing proposals that combine technical and commercial considerations.
\end{alertblock}
\end{column}
\end{columns}

\end{frame}

\begin{frame}[t]{Emotional Experience During the Task}
\small
\begin{columns}[T,onlytextwidth]
\begin{column}{.48\textwidth}
\figureplaceholder{6}{Positive emotions}{3.05cm}\begin{block}{Positive experience}
Individuals with AI show a \textbf{0.457 SD} gain in positive emotions. Teams with AI show a \textbf{0.635 SD} gain.
\end{block}
\end{column}
\begin{column}{0.48\textwidth}
\figureplaceholder{7}{Negative emotions}{3.05cm}\begin{block}{Reduced negative emotions}
The corresponding changes are \textbf{$-0.233$ SD} and \textbf{$-0.235$ SD} relative to solo work without AI.
\end{block}
\end{column}
\end{columns}

\end{frame}

\begin{frame}[t]{Idea Generation and Selection}
\small
\begin{columns}[T,onlytextwidth]
\begin{column}{.62\textwidth}
\figureplaceholder{8}{Innovation process decomposition}{4.8cm}\end{column}
\begin{column}{0.34\textwidth}
\begin{block}{Stronger initial ideas}
AI-assisted groups generate higher-quality ideas. This stronger pool leads to better selected concepts in absolute terms.
\end{block}
\begin{alertblock}{Selection still matters}
Human-only teams choose their best idea about \textbf{50\%} of the time, compared with roughly \textbf{37\%} in AI groups. Generating options and evaluating them are distinct capabilities.
\end{alertblock}
\end{column}
\end{columns}

\end{frame}

\begin{frame}[t]{Exceptional Solutions}
\small
\begin{columns}[T,onlytextwidth]
\begin{column}{.60\textwidth}
\figureplaceholder{9}{Probability of a top-decile solution}{4.7cm}
\end{column}
\begin{column}{0.36\textwidth}
\begin{block}{Finding}
Teams with AI are \textbf{9.2 percentage points} more likely than solo workers without AI to produce a top-decile solution. The baseline probability is \textbf{5.8\%}.
\end{block}
\begin{alertblock}{Interpretation}
The combination of human teamwork and AI can increase the chance of producing highly rated proposals, extending the analysis beyond average quality.
\end{alertblock}
\end{column}
\end{columns}

\end{frame}

